% ============================================================
% Paper 89 (DE): Mersenne-Primzahlen, gerade vollkommene Zahlen
%                und das Unendlichkeitsproblem
% Autor: Michael Fuhrmann
% Specialist Mathematics Project, Build 168
% Datum: 2026-03-12
% ============================================================
\documentclass[12pt,a4paper]{amsart}

\usepackage{amsmath,amssymb,amsthm,mathtools,enumitem,booktabs,array,hyperref}
\usepackage[utf8]{inputenc}
\usepackage[T1]{fontenc}
\usepackage{geometry}
\geometry{margin=2.5cm}

% ---- Theorem-Umgebungen (Deutsch) ----
\newtheorem{theorem}{Theorem}[section]
\newtheorem{lemma}[theorem]{Lemma}
\newtheorem{corollary}[theorem]{Korollar}
\newtheorem{proposition}[theorem]{Proposition}
\newtheorem{satz}[theorem]{Satz}
\newtheorem{korollar}[theorem]{Korollar}
\theoremstyle{definition}
\newtheorem{definition}[theorem]{Definition}
\newtheorem{definition_de}[theorem]{Definition}
\newtheorem{example}[theorem]{Beispiel}
\theoremstyle{remark}
\newtheorem{remark}[theorem]{Bemerkung}
\newtheorem{bemerkung}[theorem]{Bemerkung}
\newtheorem{conjecture}[theorem]{Vermutung}
\newtheorem{vermutung}[theorem]{Vermutung}

% ---- Operatoren ----
\newcommand{\N}{\mathbb{N}}
\newcommand{\Z}{\mathbb{Z}}
\newcommand{\Q}{\mathbb{Q}}
\newcommand{\R}{\mathbb{R}}
\newcommand{\F}{\mathbb{F}}

% ============================================================
\begin{document}

\title{Mersenne-Primzahlen, gerade vollkommene Zahlen\\
und das Unendlichkeitsproblem}

\author{Michael Fuhrmann}
\address{Specialist Mathematics Project, Build 168}
\email{specialist-maths@project.local}
\date{2026-03-12}

\begin{abstract}
Der Satz von Euklid und Euler stellt eine vollständige Bijektion zwischen
geraden vollkommenen Zahlen und Mersenne-Primzahlen her: Jede gerade vollkommene
Zahl hat die Form $2^{p-1}(2^p - 1)$, wobei $2^p - 1$ eine Mersenne-Primzahl ist.
Stand Oktober 2024 sind genau 52 Mersenne-Primzahlen bekannt; die größte ist
$M_{136\,279\,841} = 2^{136\,279\,841} - 1$, entdeckt am 12.\ Oktober 2024
durch Luke Durant im Rahmen des GIMPS-Projekts (\textit{Great Internet Mersenne Prime Search})
— die erste Mersenne-Primzahl mit einem Exponenten über $10^8$.
Die zentrale offene Frage lautet: Gibt es unendlich viele Mersenne-Primzahlen?
Diese Arbeit präsentiert den Satz von Euklid--Euler, den Lucas--Lehmer-Primalitätstest,
eine Übersicht bekannter Mersenne-Primzahlen, Dichteheuristiken von Lenstra und
Wagstaff sowie Verbindungen zu Kreisteilungspolynomen und Wagstaffs Repunit-Vermutung.
\end{abstract}

\maketitle

\tableofcontents

% ============================================================
\section{Einleitung}
% ============================================================

Eine \emph{Mersenne-Zahl} ist eine ganze Zahl der Form $M_n = 2^n - 1$.
Wenn $M_p = 2^p - 1$ prim ist (was voraussetzt, dass $p$ selbst prim ist),
heißt sie \emph{Mersenne-Primzahl}.
Die ersten lauten:
\[
  M_2 = 3, \quad M_3 = 7, \quad M_5 = 31, \quad M_7 = 127, \quad
  M_{13} = 8191.
\]
Mersenne-Primzahlen sind durch folgenden klassischen Satz eng mit
vollkommenen Zahlen verknüpft:

\begin{satz}[Euklid--Euler, vollständig]
\label{satz:ee}
Eine gerade positive ganze Zahl $n$ ist genau dann vollkommen, wenn
$n = 2^{p-1}(2^p - 1)$ mit einer Mersenne-Primzahl $2^p - 1$ gilt.
\end{satz}

Diese vollständige Klassifikation gerader vollkommener Zahlen --- eines der
elegantesten Ergebnisse der Zahlentheorie --- reduziert das Problem auf das
Studium von Mersenne-Primzahlen.
Die Frage nach deren Unendlichkeit gehört zu den berühmtesten offenen Problemen.

\begin{vermutung}[Unendlich viele Mersenne-Primzahlen]
\label{verm:unendlich}
Es gibt unendlich viele Mersenne-Primzahlen, also auch unendlich viele
gerade vollkommene Zahlen.
\end{vermutung}

% ============================================================
\section{Der Satz von Euklid und Euler}
% ============================================================

\subsection{Euklids Richtung}

\begin{satz}[Euklid, Elemente IX.36, $\approx$ 300 v.\,Chr.]
Ist $2^p - 1$ prim, so ist $n = 2^{p-1}(2^p - 1)$ vollkommen.
\end{satz}

\begin{proof}
Sei $q = 2^p - 1$ prim. Dann gilt:
\[
  \sigma(n) = \sigma(2^{p-1}) \cdot \sigma(q)
  = (2^p - 1)(q + 1)
  = (2^p - 1) \cdot 2^p
  = 2 \cdot 2^{p-1}(2^p - 1) = 2n. \qedhere
\]
\end{proof}

\subsection{Eulers Richtung}

\begin{satz}[Euler, $\approx$ 1747]
Jede gerade vollkommene Zahl hat die Form $2^{p-1}(2^p - 1)$
mit einer Mersenne-Primzahl $2^p - 1$.
\end{satz}

\begin{proof}
Sei $n$ eine gerade vollkommene Zahl. Schreibe $n = 2^{k-1} m$ mit ungeradem
$m$ und $k \geq 2$.
Da $\sigma$ multiplikativ ist und $\gcd(2^{k-1}, m) = 1$:
\[
  \sigma(n) = \sigma(2^{k-1})\sigma(m) = (2^k - 1)\sigma(m).
\]
Mit $\sigma(n) = 2n = 2^k m$ folgt:
\[
  (2^k - 1)\sigma(m) = 2^k m.
\]
Da $\gcd(2^k - 1, 2^k) = 1$, muss $(2^k - 1) \mid m$ gelten.
Schreibe $m = (2^k - 1) \cdot t$. Dann folgt $\sigma(m) = 2^k t = m + t$.
Also hat $m$ genau zwei Teiler: $m$ selbst und $t$.
Daher ist $m$ prim und $t = 1$, d.\,h.\ $m = 2^k - 1$ ist Mersenne-Primzahl.
\end{proof}

\begin{korollar}
Die geraden vollkommenen Zahlen stehen über $M_p \mapsto 2^{p-1} M_p$
in Bijektion zu den Mersenne-Primzahlen.
\end{korollar}

% ============================================================
\section{Notwendige Bedingung: Der Exponent muss prim sein}
% ============================================================

\begin{lemma}
\label{lem:exprim}
Ist $M_n = 2^n - 1$ prim, so ist $n$ prim.
\end{lemma}

\begin{proof}
Sei $n = ab$ mit $1 < a, b < n$. Dann gilt
$2^n - 1 = (2^a - 1)((2^a)^{b-1} + \cdots + 1)$,
also ist $2^a - 1 \mid 2^n - 1$ ein echter nichttrivialer Teiler.
Widerspruch zur Primalität.
\end{proof}

\begin{bemerkung}
Die Umkehrung von Lemma~\ref{lem:exprim} gilt nicht:
$M_{11} = 2047 = 23 \times 89$ ist zusammengesetzt, obwohl $11$ prim ist.
\end{bemerkung}

% ============================================================
\section{Der Lucas--Lehmer-Primalitätstest}
% ============================================================

Das Hauptwerkzeug zum Testen, ob $M_p = 2^p - 1$ prim ist, ist der
Lucas--Lehmer-Test, der weit effizienter als allgemeine Primalitätstests
für Zahlen dieser speziellen Form ist.

\begin{definition_de}
Die \emph{Lucas--Lehmer-Folge} ist definiert durch:
\[
  s_0 = 4, \qquad s_{k+1} = s_k^2 - 2 \pmod{M_p}.
\]
\end{definition_de}

\begin{satz}[Lucas 1878, Lehmer 1930]
\label{satz:LL}
Für eine ungerade Primzahl $p$ ist $M_p = 2^p - 1$ genau dann prim, wenn
$s_{p-2} \equiv 0 \pmod{M_p}$.
\end{satz}

\begin{proof}[Beweisidee]
Der Beweis verwendet die Theorie der Lucas-Folgen.
Betrachte den Ring $\Z[\omega]$ mit $\omega = 2 + \sqrt{3}$.
Dann gilt $s_k = \omega^{2^k} + \omega^{-2^k}$ modulo $M_p$.
Ist $M_p$ prim, so folgt mittels Fermatschen Kleinen Satzes
$\omega^{M_p+1} = \omega^{2^p} \equiv 1 \pmod{M_p}$
im Ring $\F_{M_p}[\sqrt{3}]$, was $s_{p-2} \equiv 0 \pmod{M_p}$ liefert.
Die Umkehrung nutzt die Ordnung von $\omega$ modulo $M_p$.
\end{proof}

\begin{bemerkung}
Die Zeitkomplexität des Lucas--Lehmer-Tests beträgt $O(p^2 \log p \log \log p)$
(mit schneller Multiplikation), was ihn auch für Exponenten $p$ in den
Zehnmillionen praktikabel macht. Er wird von GIMPS verwendet.
\end{bemerkung}

\begin{example}
Test für $M_5 = 31$: $s_0 = 4$, $s_1 = 14$, $s_2 = 194 \equiv 8 \pmod{31}$,
$s_3 = 62 \equiv 0 \pmod{31}$. Da $s_{5-2} = s_3 = 0$, ist $M_5 = 31$ prim. $\checkmark$
\end{example}

% ============================================================
\section{Bekannte Mersenne-Primzahlen}
% ============================================================

Stand Oktober 2024 sind genau 52 Mersenne-Primzahlen bekannt.
Ihre Exponenten $p$ lauten (ausgewählte):

\begin{center}
\begin{tabular}{rl|rl}
\toprule
Nr. & Exponent $p$ & Nr. & Exponent $p$ \\
\midrule
1 & 2 & 2 & 3 \\
3 & 5 & 4 & 7 \\
5 & 13 & 6 & 17 \\
7 & 19 & 8 & 31 \\
9 & 61 & 10 & 89 \\
\vdots & \vdots & \vdots & \vdots \\
48 & 57885161 & 49 & 74207281 \\
50 & 77232917 & 51 & 82589933 \\
52 & 136279841 & & \\
\bottomrule
\end{tabular}
\end{center}

\begin{bemerkung}
Die 51.\ Mersenne-Primzahl $M_{82\,589\,933}$ wurde am 7.\ Dezember 2018 von
Patrick Laroche im Rahmen von GIMPS entdeckt.
Sie hat $24\,862\,048$ Dezimalstellen.
\end{bemerkung}

\begin{bemerkung}[52.\ Mersenne-Primzahl, Oktober 2024]
\label{bem:52te}
Am 12.\ Oktober 2024 gab GIMPS die Entdeckung der 52.\ bekannten
Mersenne-Primzahl bekannt:
\[
    M_{136\,279\,841} = 2^{136\,279\,841} - 1.
\]
Dies ist die erste Mersenne-Primzahl mit einem Exponenten über $10^8$.
Sie hat $41\,024\,320$ Dezimalstellen und ist damit die größte bekannte
Primzahl zum Zeitpunkt dieser Arbeit.  Die Entdeckung wurde von Luke Durant
mittels eines GPU-beschleunigten GIMPS-Clients gemacht~\cite{GIMPS2024b}.
\end{bemerkung}

% ============================================================
\section{Das GIMPS-Projekt}
% ============================================================

\begin{definition_de}
GIMPS (\textit{Great Internet Mersenne Prime Search}) ist ein 1996 von George
Woltman gegründetes verteiltes Rechenprojekt, das die Testung von Mersenne-Zahlen
auf Primalität mit dem Lucas--Lehmer-Test koordiniert.
\end{definition_de}

GIMPS nutzt folgende Strategie:
\begin{enumerate}[label=(\arabic*)]
  \item Zuweisung ungetesteter Exponenten $p$ an freiwillige Computer.
  \item Durchführung des Lucas--Lehmer-Tests (mit vorheriger Probedivision).
  \item Gegenprüfung (double-check) zur Ergebnisverifizierung.
  \item Meldung verifizierter Primzahlen.
\end{enumerate}

GIMPS hat seit 1996 insgesamt 17 der größten bekannten Mersenne-Primzahlen
entdeckt. Das Projekt gewann den EFF-Preis (Electronic Frontier Foundation)
von \$100\,000 für die erste Primzahl mit über 10 Millionen Stellen
($M_{43\,112\,609}$, 2008).
Der \$150\,000-Preis für die erste 100-Millionen-stellige Primzahl ist
Stand 2024 noch offen: Die größte bekannte Mersenne-Primzahl $M_{136\,279\,841}$
hat circa $41$ Millionen Stellen.

% ============================================================
\section{Dichteheuristiken: Die Lenstra--Wagstaff-Vermutung}
% ============================================================

\begin{vermutung}[Lenstra--Wagstaff]
\label{verm:LW}
Die Anzahl der Mersenne-Primzahlen $M_p = 2^p - 1$ mit $p \leq x$ ist
asymptotisch
\[
  \sim \frac{e^\gamma}{\log 2} \cdot \log\log x
  \approx 2{,}5695 \cdot \log\log x,
\]
wobei $\gamma \approx 0{,}5772$ die Euler--Mascheroni-Konstante ist.
\end{vermutung}

\begin{bemerkung}
Diese Heuristik entsteht durch probabilistische Modellierung der Primalität
von $M_p$.
Nach dem Primzahlsatz hat $M_p \approx 2^p$ etwa $p \log 2$ Stellen, also ist
die ``Wahrscheinlichkeit'' der Primalität ungefähr $1/(p \log 2)$.
Summation über Primzahlen $p \leq x$ liefert:
\[
  \sum_{\substack{p \leq x \\ p \text{ prim}}} \frac{1}{p \log 2}
  \approx \frac{1}{\log 2} \log\log x.
\]
Der Faktor $e^\gamma$ entsteht durch sorgfältigere Analyse der Spezialstruktur
der Mersenne-Zahlen.
\end{bemerkung}

\begin{bemerkung}
Vermutung~\ref{verm:LW} impliziert die Existenz unendlich vieler Mersenne-Primzahlen
(da $\log\log x \to \infty$), konsistent mit Vermutung~\ref{verm:unendlich}.
Sie ersetzt jedoch keinen Beweis.
\end{bemerkung}

% ============================================================
\section{Kreisteilungspolynome und Mersenne-Zahlen}
% ============================================================

\begin{definition_de}
Das \emph{$n$-te Kreisteilungspolynom} ist
\[
  \Phi_n(x) = \prod_{\substack{k=1 \\ \gcd(k,n)=1}}^n (x - e^{2\pi i k/n}).
\]
\end{definition_de}

Mersenne-Zahlen erfüllen die Faktorisierung:
\[
  2^n - 1 = \prod_{d \mid n} \Phi_d(2).
\]

\begin{proposition}
Jeder Primteiler $q$ von $\Phi_p(2) = 2^{p-1} + \cdots + 1$ erfüllt
$q \equiv 1 \pmod{p}$ für Primzahlen $p$.
\end{proposition}

\begin{proof}
Ein Primteiler $q \mid \Phi_p(2)$ hat die Ordnung $\ord_q(2) = p$.
Nach dem Fermatschen Kleinen Satz gilt $p \mid q-1$, also $q \equiv 1 \pmod{p}$.
\end{proof}

\begin{korollar}
Jeder Primteiler von $M_p = 2^p - 1$ ist $\equiv 1 \pmod{p}$
(für ungerade Primzahlen $p$).
Dies schränkt die Kandidaten stark ein und wird in der Probedivision ausgenutzt.
\end{korollar}

% ============================================================
\section{Wagstaffs Repunit-Verbindung}
% ============================================================

\begin{definition_de}
Eine \emph{Mersenne-Zahl zur Basis $b$} ist $M_{p,b} = (b^p - 1)/(b - 1)$
für $b \geq 2$.
Für $b = 10$: die Repunits $R_p = \underbrace{11\cdots1}_p$.
\end{definition_de}

\begin{vermutung}[Wagstaff 1983]
Die Anzahl der Primzahlen der Form $(b^p - 1)/(b-1)$ mit Primzahl $p \leq x$
ist für jedes feste $b \geq 2$ asymptotisch $\sim (e^\gamma / \log b) \log\log x$.
\end{vermutung}

\begin{bemerkung}
Diese Vermutung vereinheitlicht die Mersenne-Primzahl-Heuristik für verschiedene
Basen und unterstützt die Erwartung, dass für jede Basis unendlich viele solche
Primzahlen existieren.
\end{bemerkung}

% ============================================================
\section{Eigenschaften gerader vollkommener Zahlen}
% ============================================================

\begin{proposition}
Jede gerade vollkommene Zahl $n = 2^{p-1}(2^p-1)$ erfüllt:
\begin{enumerate}[label=(\roman*)]
  \item $n$ ist eine Dreieckszahl: $n = 1 + 2 + \cdots + (2^p - 1)$;
  \item Die Binärdarstellung von $n$ besteht aus $p$ Einsen gefolgt von
        $p-1$ Nullen: $\underbrace{1\cdots1}_{p}\underbrace{0\cdots0}_{p-1}$;
  \item $n \equiv 1 \pmod{9}$ für $p > 3$;
  \item $n$ ist nicht durch $3$ teilbar (für $p > 2$).
\end{enumerate}
\end{proposition}

\begin{proof}[Beweis von (i)]
$\sum_{k=1}^{2^p-1} k = (2^p-1) \cdot 2^p / 2 = 2^{p-1}(2^p-1) = n$.
\end{proof}

% ============================================================
\section{Primitive Primteiler und Zsygmondy-Satz}
% ============================================================

\begin{proposition}[Carmichael für Mersenne-Zahlen]
Für $n > 2$ besitzt $M_n = 2^n - 1$ einen Primteiler, der $M_k$ für kein
$1 \leq k < n$ teilt (einen \emph{primitiven Primteiler}).
\end{proposition}

\begin{proof}
Dies folgt aus dem Satz von Zsygmondy (Birkhoff--Vandiver), angewandt auf
$a = 2$, $b = 1$, $n \geq 3$.
\end{proof}

\begin{bemerkung}
Der Zsygmondy-Satz zeigt, dass die Mersenne-Zahlen immer neue Primteiler
einbringen, was ein schwaches Argument für die Unendlichkeit von
Mersenne-Primzahlen liefert --- obwohl neue Primteiler nicht neue
Mersenne-Primzahlen bedeuten.
\end{bemerkung}

% ============================================================
\section{Das Unendlichkeitsproblem}
% ============================================================

\begin{vermutung}[Wiederholt: Unendlich viele Mersenne-Primzahlen]
Es gibt unendlich viele Primzahlen der Form $2^p - 1$.
\end{vermutung}

Trotz umfangreicher heuristischer Belege und rechnerischer Unterstützung ist
kein Beweis der Unendlichkeit bekannt.
Die Hauptschwierigkeit liegt darin, dass die Primzahlen $p$ in $M_p = 2^p - 1$
gleichzeitig erfüllen müssen:
\begin{itemize}[nosep]
  \item $p$ ist selbst prim;
  \item $M_p$ besteht den Lucas--Lehmer-Test.
\end{itemize}

\begin{bemerkung}[Warum es schwer ist]
Die Unendlichkeit von Primzahlen der Form $f(n)$ für Polynome $f$ --- selbst
für lineare (Dirichletscher Primzahlsatz) --- erfordert tiefe analytische Methoden.
Für Exponentialfunktionen wie $2^p - 1$ fehlen der analytischen Zahlentheorie
derzeit geeignete Werkzeuge für direkte Unendlichkeitsbeweise.
Siebmethoden sind ineffektiv, weil die Folge zu schnell wächst.
\end{bemerkung}

\begin{bemerkung}[Keine obere Schranke]
Es ist ebenso unbekannt, ob die Menge der Mersenne-Primzahlen \emph{endlich}
ist. Dies zu beweisen würde erfordern, dass der Lucas--Lehmer-Test ab einem
gewissen Punkt immer versagt --- was genauso unlösbar erscheint.
\end{bemerkung}

% ============================================================
\section{Verbindungen zu weiteren Strukturen}
% ============================================================

\subsection{Amikable Zahlen}

Zwei Zahlen $m, n$ heißen \emph{amikabel}, wenn $s(m) = n$ und $s(n) = m$
gilt. Euler fand 58 amikable Paare und vermutete deren Unendlichkeit;
auch dies ist offen.

\subsection{Wieferich-Primzahlen}

Eine Primzahl $p$ heißt \emph{Wieferich-Primzahl}, wenn
$2^{p-1} \equiv 1 \pmod{p^2}$ gilt.
Nur zwei sind bekannt: $1093$ und $3511$.
Die Verbindung zu Mersenne-Primzahlen ist heuristisch (keine Mersenne-Primzahl
sollte eine Wieferich-Primzahl als Exponent haben), aber ohne Beweis.

\subsection{Doppelt vollkommene Zahlen}

Gerade vollkommene Zahlen erfüllen $\sigma(n) = 2n$ und sind der einzige
bekannte vollständig klassifizierte Typ vollkommener Zahlen.

% ============================================================
\section{Schluss}
% ============================================================

Die Theorie der Mersenne-Primzahlen und gerader vollkommener Zahlen verbindet
die antike Zahlentheorie Euklids mit modernster Computermathematik (GIMPS,
Lucas--Lehmer).
Der Satz von Euklid--Euler liefert eine vollständige Klassifikation gerader
vollkommener Zahlen und reduziert das Problem auf das Studium von Primzahlen
der Form $2^p - 1$.

Die zentralen offenen Fragen bleiben:
\begin{enumerate}[label=(\arabic*)]
  \item Gibt es unendlich viele Mersenne-Primzahlen?
  \item Gibt es ungerade vollkommene Zahlen? (Siehe Paper~88)
  \item Kann die Lenstra--Wagstaff-Dichtevermutung bewiesen werden?
  \item Gibt es ein Muster unter den Exponenten der Mersenne-Primzahlen?
\end{enumerate}

Diese scheinbar einfach zu formulierenden Fragen berühren einige der tiefsten
ungelösten Probleme der analytischen und rechnerischen Zahlentheorie.

% ============================================================
\begin{thebibliography}{99}

\bibitem{Euklid300}
Euklid,
\textit{Elemente},
Buch IX, Satz 36, $\approx$ 300 v.\,Chr.

\bibitem{Euler1747}
L.\ Euler,
\textit{De numeris amicabilibus},
Nova Acta Eruditorum (1747).

\bibitem{Lucas1878}
E.\ Lucas,
\textit{Théorie des fonctions numériques simplement périodiques},
Amer.\ J.\ Math.\ \textbf{1} (1878), 184--240, 289--321.

\bibitem{Lehmer1930}
D.\ H.\ Lehmer,
\textit{An extended theory of Lucas' functions},
Ann.\ Math.\ \textbf{31} (1930), 419--448.

\bibitem{Wagstaff1983}
S.\ S.\ Wagstaff Jr.,
\textit{Divisors of Mersenne numbers},
Math.\ Comp.\ \textbf{40} (1983), 385--397.

\bibitem{Caldwell2024}
C.\ K.\ Caldwell,
\textit{Mersenne Primes: History, Theorems and Lists},
The Prime Pages, \url{https://primes.utm.edu/mersenne/}, 2024.

\bibitem{GIMPS2024}
GIMPS,
\textit{Great Internet Mersenne Prime Search},
\url{https://www.mersenne.org/}, 2024.

\bibitem{Laroche2018}
P.\ Laroche,
\textit{Announcement of $M_{82\,589\,933}$},
GIMPS-Pressemitteilung, Dezember 2018.

\bibitem{Zsygmondy1892}
K.\ Zsygmondy,
\textit{Zur Theorie der Potenzreste},
Monatsh.\ Math.\ Phys.\ \textbf{3} (1892), 265--284.

\bibitem{CrandallPomerance2005}
R.\ Crandall und C.\ Pomerance,
\textit{Prime Numbers: A Computational Perspective},
2.\ Aufl., Springer, 2005.

\bibitem{GIMPS2024b}
GIMPS,
\textit{GIMPS entdeckt größte bekannte Primzahl: $2^{136\,279\,841}-1$},
Pressemitteilung, 12.\ Oktober 2024,
\url{https://www.mersenne.org/}.

\bibitem{Ribenboim1996}
P.\ Ribenboim,
\textit{The New Book of Prime Number Records},
Springer, New York, 1996.

\bibitem{HardyWright2008}
G.\ H.\ Hardy und E.\ M.\ Wright,
\textit{An Introduction to the Theory of Numbers},
6.\ Aufl., Oxford University Press, 2008.

\end{thebibliography}

\end{document}
